% =====================================================================
%  Pacchetti
%  La classe Sapthesis carica già: geometry, fontenc (T1), lmodern,
%  caption (etichetta piccola e in grassetto, separatore "."), graphicx,
%  color, booktabs, amsmath, etoolbox, fancyhdr.
% =====================================================================

% --- Lingua e testo ----------------------------------------------------
\usepackage[italian]{babel}
\usepackage[autostyle,italian=guillemets]{csquotes}
\usepackage{microtype}
\usepackage{icomma}        % virgola decimale in modo matematico: $0,604$
\usepackage{enumitem}
\setlist{itemsep=0.15em, parsep=0.15em, topsep=0.3em}

% Paragrafi separati da una riga vuota e senza rientro, come nella tesi del 2022.
\usepackage[skip=0.75\baselineskip plus 2pt, indent=0pt]{parskip}

% --- Colori (xcolor sostituisce color, caricato dalla classe) ------------
\usepackage[table,dvipsnames]{xcolor}
% =====================================================================
%  Palette colori
%  Ricavata dal PDF della tesi del 2022 (legacy/tesi-magistrale/), leggendo
%  i valori nei content stream e nelle annotazioni dei link:
%    - rosso Sapienza RGB(130,36,51): frontespizio (definito dalla classe
%      Sapthesis come `sapred', ridefinito qui per usarlo già nel preambolo);
%    - cornici dei link di hyperref: rosso [1 0 0] per i riferimenti interni,
%      verde [0 1 0] per le citazioni, ciano [0 1 1] per gli URL.
%  Tutti gli altri colori sono tinte di questi, oppure i colori ufficiali
%  delle classi ISPRS, che servono a leggere maschere e legende.
% =====================================================================

% --- Colori di identità (identici alla tesi del 2022) -----------------
\definecolor{SapRosso}{rgb}{0.5098039,0.1411765,0.2}   % = sapred, #822433
\definecolor{LinkRosso}{rgb}{1,0,0}                    % riferimenti interni
\definecolor{CitaVerde}{rgb}{0,1,0}                    % citazioni
\definecolor{UrlCiano}{rgb}{0,1,1}                     % URL

% --- Tinte di supporto per schemi e tabelle ---------------------------
\colorlet{SapRossoScuro}{SapRosso!80!black}
\colorlet{SapRossoMedio}{SapRosso!60}
\colorlet{SapRossoTenue}{SapRosso!30}
\colorlet{SapRossoChiaro}{SapRosso!12}
\colorlet{SapRossoVelo}{SapRosso!5}
\colorlet{SapGrigio}{black!55}
\colorlet{SapGrigioTenue}{black!25}
\colorlet{SapGrigioChiaro}{black!8}

% --- Classi ISPRS (palette ufficiale delle maschere di Potsdam/Vaihingen)
\definecolor{ClsImpervious}{RGB}{255,255,255}   % superfici impermeabili
\definecolor{ClsBuilding}{RGB}{0,0,255}         % edifici
\definecolor{ClsLowVeg}{RGB}{0,255,255}         % vegetazione bassa
\definecolor{ClsTree}{RGB}{0,255,0}             % alberi
\definecolor{ClsCar}{RGB}{255,255,0}            % auto
\definecolor{ClsClutter}{RGB}{255,0,0}          % clutter / sfondo

% --- Ruoli semantici negli schemi --------------------------------------
\colorlet{ColEncoder}{SapRossoChiaro}     % encoder / backbone
\colorlet{ColDecoder}{SapGrigioChiaro}    % decoder
\colorlet{ColAttn}{SapRossoTenue}         % blocchi di attenzione
\colorlet{ColQuant}{SapRosso}             % quantizzazione vettoriale
\colorlet{ColFreccia}{SapGrigio}          % frecce di flusso
\colorlet{ColGradiente}{SapRosso}         % percorsi del gradiente

% --- Split dei dati -----------------------------------------------------
\colorlet{SplitTrain}{SapRossoMedio}
\colorlet{SplitVal}{SapRossoTenue}
\colorlet{SplitTest}{SapGrigio}
\colorlet{SplitEscl}{white}

% --- Righe evidenziate e segnaposto -------------------------------------
\colorlet{RigaRiferimento}{SapRossoVelo}  % righe di riferimento nelle tabelle
\colorlet{ColTodo}{SapRossoChiaro}        % sfondo delle note \todo
\colorlet{ColTodoBordo}{SapRosso}


% --- Matematica ----------------------------------------------------------
\usepackage{amssymb}
\usepackage{amsthm}
\usepackage{mathtools}
\usepackage{bm}

% --- Tabelle -------------------------------------------------------------
\usepackage{array}
\usepackage{multirow}
\usepackage{makecell}
\usepackage{tabularx}
\usepackage{siunitx}
\sisetup{output-decimal-marker={,}, group-separator={.}, group-minimum-digits=4,
         detect-weight=true, detect-family=true}
\renewcommand{\arraystretch}{1.15}

% --- Figure e schemi -----------------------------------------------------
\usepackage{float}          % [H]: posizione fissa (usata nei capitoli con segnaposto)
\usepackage{subcaption}
\usepackage{tikz}
\usetikzlibrary{positioning, arrows.meta, calc, fit, backgrounds,
                shapes.geometric, shapes.misc, matrix, decorations.pathreplacing,
                patterns, shadows}
\usepackage{pgfplots}
\pgfplotsset{compat=1.18}

% --- Algoritmi -----------------------------------------------------------
\usepackage{algorithm}
\usepackage{algpseudocode}
\floatname{algorithm}{Algoritmo}
\renewcommand{\listalgorithmname}{Elenco degli algoritmi}
\algrenewcommand\algorithmicrequire{\textbf{Input:}}
\algrenewcommand\algorithmicensure{\textbf{Output:}}
\algrenewcommand\algorithmicfor{\textbf{per}}
\algrenewcommand\algorithmicdo{\textbf{esegui}}
\algrenewcommand\algorithmicend{\textbf{fine}}
\algrenewcommand\algorithmicif{\textbf{se}}
\algrenewcommand\algorithmicthen{\textbf{allora}}
\algrenewcommand\algorithmicreturn{\textbf{restituisci}}
\algnewcommand\algorithmicforeach{\textbf{per ogni}}
\algdef{S}[FOR]{ForEach}[1]{\algorithmicforeach\ #1\ \algorithmicdo}

% --- Note di lavoro e segnaposto -----------------------------------------
% Per la versione finale: \bozzafalse in main.tex (le note spariscono).
\ifbozza
  \usepackage[colorinlistoftodos, italian, textsize=footnotesize,
              color=ColTodo, bordercolor=ColTodoBordo, linecolor=ColTodoBordo]{todonotes}
\else
  \usepackage[disable]{todonotes}
\fi
\setlength{\marginparwidth}{2.2cm}

% --- Collegamenti (stesse cornici colorate della tesi del 2022) -----------
\usepackage{hyperref}
\hypersetup{
  pdftitle    = {\TitoloTesi},
  pdfauthor   = {Matteo Campana},
  pdfsubject  = {Tesi di Laurea Magistrale in Engineering in Computer Science},
  pdfkeywords = {segmentazione semantica, telerilevamento, attenzione lineare,
                 quantizzazione vettoriale, apprendimento auto-supervisionato, DINO},
  colorlinks      = false,
  pdfborder       = {0 0 1},
  linkbordercolor = LinkRosso,
  citebordercolor = CitaVerde,
  urlbordercolor  = UrlCiano,
  bookmarksnumbered = true,
}
\usepackage[italian, noabbrev]{cleveref}
\crefname{algorithm}{algoritmo}{algoritmi}
\Crefname{algorithm}{Algoritmo}{Algoritmi}
\crefname{ipotesi}{ipotesi}{ipotesi}
\Crefname{ipotesi}{Ipotesi}{Ipotesi}
